\documentclass{article}
\usepackage[zh]{homework}

\config{代数\,-\,0}{2026 春季}{8}

\begin{document}


\begin{problem}[6A-14]
设 $v \in V$($v \ne 0$). 证明: $v/\lVert v \rVert$ 是 $V$ 的单位球面上唯一最接近 $v$ 的元素. 更准确地说, 证明: 如果 $u \in V$ 且 $\lVert u \rVert = 1$, 那么
\[
\left\lVert v - \frac{v}{\lVert v \rVert} \right\rVert \le \lVert v - u \rVert,
\]
且仅当 $u = v/\lVert v \rVert$ 时取等.
\end{problem}

\begin{proof}
  注意到
  \begin{align*}
    \abs{v-u}^2
    &=\inprod{v-u}{v-u}
    =\abs{v}^2 - 2\Re\inprod{v}{u} + \abs{u}^2
    \ge\abs{v}^2-2\abs{v}+1\\
    &=\left( \abs{v}-1 \right)^2=\abs{v-\frac{v}{\abs{v}}}^2.
  \end{align*}
  其中用到 $\abs{u}=1$. 由柯西不等式, $\Re\inprod{v}{u} \le \abs{\inprod{v}{u}} \le \abs{v}\abs{u}=\abs{v}$, 因此不等式成立. 若取等, 则 $\abs{\inprod{v}{u}}=\abs{v}$ 且 $\Re\inprod{v}{u}=\abs{v}$, 从而 $\inprod{v}{u}=\abs{v}$ 并有 $u=v/\abs{v}$. 反之当 $u=v/\abs{v}$ 时等号成立.
\end{proof}



\begin{problem}[6A-18]
\textbf{(a)} 设 $f : [1, \infty) \to [0, \infty)$ 连续. 证明:
\[
\left(\int_1^{\infty} f \right)^2 \le \int_1^{\infty} x^2 (f(x))^2\,\d x.
\]
\textbf{(b)} 当连续函数 $f : [1, \infty) \to [0, \infty)$ 是什么函数时, 以上不等式取等且两边均有限?
\end{problem}

\begin{solution}
  (a) 注意到 \( \int_1^\infty x^{-2}\d x=1 \). 由柯西不等式, 
  \[ \left( \int_{1}^{\infty}x^2f^2(x)\d x \right)\left( \int_{1}^{\infty} x^{-2}\d x \right) \ge \left( \int_{1}^{\infty} f(x) \d x \right)^2 \]
  故不等式成立.\par
  (b) 当且仅当 \( f(x)=\lambda x^{-2} \) 对某个 \( \lambda \) 成立时, 不等式取等. 此时两边均为 \( \lambda^2 \).
\end{solution}



\begin{problem}[6A-22]
证明: 如果 $u, v \in V$, 那么
\[
\lVert u + v \rVert\,\lVert u - v \rVert \le \lVert u \rVert^2 + \lVert v \rVert^2.
\]
\end{problem}

\begin{proof}
  \begin{align*}
    \abs{u+v}^2\abs{u-v}^2
    &=\inprod{u+v}{u+v}\inprod{u-v}{u-v}
    =\left( \abs{u}^2+\abs{v}^2-2\Re\inprod{u}{v} \right)\left( \abs{u}^2+\abs{v}^2+2\Re\inprod{u}{v} \right)\\
    &=(\abs{u}^2+\abs{v}^2)^2-4(\Re\inprod{u}{v})^2
    \le(\abs{u}^2+\abs{v}^2)^2.
  \end{align*}
  两边非负, 取平方根得 $\abs{u+v}\abs{u-v} \le \abs{u}^2+\abs{v}^2$, 因此原不等式成立.
\end{proof}



\begin{problem}[6A-27]
设 $V$ 是复内积空间, 证明:
\[
\langle u, v \rangle
= \frac{\lVert u + v \rVert^2 - \lVert u - v \rVert^2 + \i\,\lVert u + \i v \rVert^2 - \i\,\lVert u - \i v \rVert^2}{4}
\]
对所有 $u, v \in V$ 成立.
\end{problem}

\begin{proof}
  由内积的线性与共轭线性, 有
  \begin{align*}
    \abs{u+v}^2 &= \inprod{u}{u}+\inprod{u}{v}+\inprod{v}{u}+\inprod{v}{v},\\
    \abs{u-v}^2 &= \inprod{u}{u}-\inprod{u}{v}-\inprod{v}{u}+\inprod{v}{v},\\
    \abs{u+\i v}^2 &= \inprod{u}{u}+\inprod{v}{v}+\i\big(\inprod{v}{u}-\inprod{u}{v}\big),\\
    \abs{u-\i v}^2 &= \inprod{u}{u}+\inprod{v}{v}-\i\big(\inprod{v}{u}-\inprod{u}{v}\big).
  \end{align*}
  于是
  \[
    \abs{u+v}^2-\abs{u-v}^2=2\big(\inprod{u}{v}+\inprod{v}{u}\big),\quad
    \abs{u+\i v}^2-\abs{u-\i v}^2=2\i\big(\inprod{v}{u}-\inprod{u}{v}\big).
  \]
  代回即得结论.
\end{proof}



\begin{problem}[6A-28]
对于向量空间 $U$, 其上的范数是这样一个函数
\[
\lVert \cdot \rVert : U \to [0, \infty),
\]
其满足下列性质: $\lVert u \rVert = 0$ 当且仅当 $u = 0$; $\lVert \alpha u \rVert = \abs{\alpha}\,\lVert u \rVert$ 对所有 $\alpha \in \F$ 和所有 $u \in U$ 成立;
$\lVert u + v \rVert \le \lVert u \rVert + \lVert v \rVert$ 对所有 $u, v \in U$ 成立.
证明: 满足平行四边形等式的范数来自内积——换句话说, 要证明: 如果 $\lVert \cdot \rVert$ 是 $U$ 上满足平行四边形等式的范数, 那么 $U$ 上就存在一内积 $\langle \cdot, \cdot \rangle$, 使得对所有 $u \in U$ 有 $\lVert u \rVert = \langle u, u \rangle^{1/2}$.
\end{problem}

\begin{proof}
  假设范数 \( \abs{\cdot} \) 满足平行四边形等式:
  \[ \abs{u+v}^2+\abs{u-v}^2=2\left( \abs{u}^2+\abs{v}^2 \right). \]
  则定义内积
  \[
\langle u, v \rangle
= \frac{\abs{u+v}^2 - \abs{u-v}^2 + \i\,\abs{u+\i v}^2 - \i\,\abs{u -\i v}^2}{4}
\]\par
  下面来验证此定义满足内积的所有性质.
  \begin{itemize}
    \item \textbf{正定性. }由定义, \( \langle u, u \rangle=\abs{u}^2 \). 故对于任意的 \( u \) 均有 \( \langle u, u \rangle \ge 0 \), 且仅当 \( u=0 \) 时取等.
    \item \textbf{共轭对称性. }由定义直接计算, 并注意 \(\overline{\abs{w}^2}=\abs{w}^2\), 得 \( \overline{\langle u, v \rangle}=\langle v, u \rangle \).
    \item \textbf{可加性. }对于任意的 \( u_1,u_2,v\in U \) , 由平行四边形等式, 均有
    \begin{align*}
      &\abs{u_1+u_2+v}^2-\abs{u_1+u_2-v}^2\\
      =&2\left( \abs{u_1+v}^2+\abs{u_2}^2 \right)-\abs{u_1-u_2+v}^2-2\left( \abs{u_1-v}^2+\abs{u_2}^2 \right)+\abs{u_2-u_1+v}^2.
    \end{align*}
    交换 \( u_1 \) 与 \( u_2 \), 再与上面的式子相加, 就有
    \begin{align*}
      &\abs{u_1+u_2+v}^2-\abs{u_1+u_2-v}^2\\
      =&\abs{u_1+v}^2-\abs{u_1-v}^2+\abs{u_2+v}^2-\abs{u_2-v}^2.
    \end{align*}\par
    同理, 将 \( v \) 替换为 \( \i v \), 可以得到类似的等式. 两者结合, 就可以证明 \( \inprod{u_1+u_2}{v}=\inprod{u_1}{v}+\inprod{u_2}{v} \).
    \item \textbf{齐次性. }先来证明: 对于任意的 \( u,v\in V \) 和 \( \lambda\in\R \), 有 \( \inprod{\lambda u}{v}=\lambda\inprod{u}{v} \).\par
    首先, 由刚刚证明的可加性, 上面的命题对于任意的 \( \lambda\in\Q \) 均成立. 再由三角不等式, 知对于任意的正实数 \( \varepsilon \) 以及 \( x,y\in V \),有
    \[ \abs{x+\varepsilon y}\in [\abs{x}-\varepsilon\abs{y}, \abs{x}+\varepsilon\abs{y}], \]
    因此范数 \( \abs{\cdot} \) 是连续函数, 因而对于任意的实数 \( \lambda \), 都有 \( \inprod{\lambda u}{v}=\lambda\inprod{u}{v} \).\par
    又有
    \[
      \inprod{\i u}{v}
      = \frac{\abs{u-\i v}^2-\abs{u+\i v}^2+\i\abs{u+v}^2-\i\abs{u-v}^2}{4}
      = \i\inprod{u}{v},
    \]
    从而对任意 \( \alpha=a+b\i \) 有 \( \inprod{\alpha u}{v}=\alpha\inprod{u}{v} \). 由共轭对称性, 内积对第二变量共轭线性.
     
  \end{itemize}
  因此这是一个内积, 且 \( \lVert u \rVert=\langle u, u \rangle^{1/2} \).
\end{proof}



\begin{problem}[6A-35]
取定一个正整数 $n$. $\R^n$ 上的二阶可微实值函数 $p$ 的拉普拉斯算子(Laplacian)$\Delta p$ 是 $\R^n$ 上的函数, 定义为
\[
\Delta p = \frac{\partial^2 p}{\partial x_1^2} + \cdots + \frac{\partial^2 p}{\partial x_n^2}.
\]
如果 $\Delta p = 0$, 则称函数 $p$ 为调和的(harmonic).
$\R^n$ 上的多项式是形如 $x_1^{m_1} \cdots x_n^{m_n}$ 的函数的线性组合(其系数在 $\R$ 中), 其中 $m_1, \ldots, m_n$ 是非负整数.
设 $q$ 是 $\R^n$ 上的多项式, 证明: $\R^n$ 上存在一调和多项式 $p$, 使得 $p(x) = q(x)$ 对每个满足 $\lVert x \rVert = 1$ 的 $x \in \R^n$ 成立.
\end{problem}

\begin{proof}
  记 \( \R^n \) 上的所有多项式函数组成的线性空间为 \( V \). 考虑算子 \( T \), 即
  \[ (Tp)(x):=\Delta((1-\abs{x}^2)p(x)). \]\par
  不难验证 \( T \) 是 \( V \) 上的线性算子. 下面来证明 \( T \) 是满射. 只需要证明 \( T \) 是单射.  假设 \( p\in V \) 满足 \( Tp=0 \). 记 \( q=(1-\abs{x}^2)p(x) \), 则 \( \Delta q=0 \). 则 \( q \) 是光滑的调和函数, 并且在单位球上均为 \( 0 \). 由调和函数的经典结论, \( q \) 在单位球内也为 \( 0 \). 结合 \( q \) 为多项式, 知 \( q \) 为零多项式, 因此 \( p \) 也为零多项式. 因而 \( T \) 是单射.\par
  因此, 对于 \( R^n \) 上的任意一个多项式 \( q \), 存在多项式 \( p \) 使得
  \[ Tp=\Delta q\implies\Delta(q(x)-(1-\abs{x}^2)p(x))=0. \]\par
  因此多项式函数 \( q(x)-(1-\abs{x}^2p(x)) \) 是调和函数, 并且在单位球面上与 \( q \) 的取值相同. 这就证明了题目要求的结论.
\end{proof}



\begin{problem}[6B-9]
设 $e_1, \ldots, e_m$ 是对 $V$ 中的线性无关组 $v_1, \ldots, v_m$ 应用格拉姆-施密特过程得到的结果. 证明: $\langle v_k, e_k \rangle > 0$ 对任一 $k = 1, \ldots, m$ 成立.
\end{problem}

\begin{proof}
  由 Gram-Schmidt 过程, 有
  \[ e_k= v_k - \sum_{i=1}^{k-1} \frac{\inprod{v_k}{e_i}}{\abs{e_i}^2} e_i\implies \inprod{v_k}{e_k}=\abs{e_k}^2+ \sum_{i=1}^{k-1} \frac{\inprod{v_k}{e_i}}{\abs{e_i}^2} \inprod{e_k}{e_i}=\abs{e_k}^2>0. \]
\end{proof}



\begin{problem}[6B-21]
设 $\F = \C$, $V$ 有限维, $T \in \mathcal{L}(V)$, 且 $T$ 的所有特征值绝对值都小于 $1$. 令 $\varepsilon > 0$, 证明: 存在一正整数 $m$ 使得 $\lVert T^m v \rVert \le \varepsilon \lVert v \rVert$ 对任一 $v \in V$ 成立.
\end{problem}

\begin{proof}
  将命题加强为: 存在常数 $c > 0$ 使得对于任意的 $v \in V$ 和 $k \in \mathbb{N}$, 都有 $\abs{T^k v} \le k^nc\Lambda^k \abs{v}$, 其中 \( \Lambda \) 是 \( T \) 的所有特征值绝对值的最大值, \( n=\dim V \). \par
  对 \( n \) 归纳. \( n=1 \) 时, \( T \) 的特征值 \( \lambda \) 满足 \( \abs{\lambda}<1 \). 取 \( m \) 使得 \( \abs{\lambda}^m\le\varepsilon \), 则对于任意的 \( v\in V \), 都有
  \[ \abs{T^mv}=\abs{\lambda^m}\abs{v}\le\varepsilon\abs{v}. \]\par
  下面假设命题对小于 \( n \) 的维数成立. 设 \( \lambda \) 是 \( T \) 的一个特征值, 以及 \( e \) 是对应的特征单位向量. 记 \( W=\sspan \{e\} \), 并且将 \( \{e\} \) 扩展为 \( V \) 的一组标准正交基 \( \mathcal{B}=\{e\}\cup \mathcal{B'} \). \par
  记 \( M=\mathcal{M}(T,\mathcal{B}) \), 而 \( M'=\mathcal{M}(T/W,\mathcal{B'}+W) \), 其中 \( \mathcal{B'}+W \) 表示 \( \{v+W:v\in\mathcal{B'}\} \), 这是 \( V/W \) 的一组基. 则
  \[ M=\begin{pmatrix}
    \lambda & A\\
    0 & M'
  \end{pmatrix}, \]
  其中 \( A\in\C^{1\times(n-1)} \). 设 \( \Lambda \) 是 \( M \) 的所有特征值绝对值的最大值, 而 \( \Lambda' \) 是 \( M' \) 的所有特征值绝对值的最大值. 由于 \( M' \) 的特征值也是 \( M \) 的特征值, 因此 \( \Lambda' \le \Lambda < 1 \). \par
  由归纳假设, 存在常数 \( c' > 0 \) 使得对于任意的 \( w\in W^{\perp} \) 和 \( k\in\N \), 都有(其中 \( T' \) 是 \( M' \) 在 \( W^{\perp} \) 上对应的线性变换; 具体地, \( T'=P_{W^\perp}\circ T\circ P_{W^\perp} \))
  \[ \abs{T'^k w} \le c'k^{n-1} \Lambda'^k \abs{w}\le c'k^{n-1}\Lambda^k\abs{w}. \]\par
  记 \( a=\abs{A\transpose} \). 则对于任意的 \( v\in\C^{(n-1)\times 1} \), 都有
  \[ \abs{Av}=\abs{\inprod{v}{A\transpose}}\le a\abs{v}. \]\par
  我们来证明: 存在常数 \( c \), 使得对于任意的 \( v\in V \) 和 \( k\in\N \), 都有 \( \abs{T^k v}\le c\Lambda^k\abs{v} \). 只需要证明: 对于任意的
  \[ v=\begin{pmatrix}
    x\\
    w
  \end{pmatrix}\in\C^{n\times 1}, \text{ 其中 } x\in\C, w\in\C^{(n-1)\times 1}, \]
  均有 \( \abs{M^kv} \le c\Lambda^k\abs{v} \).\par
  由 \( M \) 的定义, 有
  \[ M^kv=\begin{pmatrix}
    \lambda^k & \sum_{i=1}^{k}\lambda^{k-i}A M'^{i-1}\\
    0 & M'^k
  \end{pmatrix}\begin{pmatrix}
    x\\
    w
  \end{pmatrix}
  =\begin{pmatrix}
    \lambda^k x + \sum_{i=1}^{k}\lambda^{k-i}A M'^{i-1}w\\
    M'^kw
  \end{pmatrix}. \]
  注意到
  \begin{align*}
    \abs{\lambda^kx+\sum_{i=1}^{k}\lambda^{k-i}A M'^{i-1}w}
    &\le \Lambda^k\abs{x} + \sum_{i=1}^{k}\Lambda^{k-i}a\abs{M'^{i-1}w}
    \le \Lambda^k\abs{x} + a c' \sum_{i=1}^{k}\Lambda^{k-i}\Lambda^{i-1}i^{n-1}\abs{w}\\
    &\le \left( 1+ac'k^n \right)\Lambda^k\abs{v}\le(1+ac')k^n\Lambda^k\abs{v}.
  \end{align*}
  再结合 \( \abs{M'^kw}\le c'k^{n-1}\Lambda^k\abs{w}\le c'k^n\Lambda^k\abs{v} \), 就有
  \[ \abs{M^kv}\le\sqrt{(1+ac')^2+(c')^2}k^n\Lambda^k\abs{v}. \]
  因此, 取 \( c=\sqrt{(1+ac')^2+(c')^2} \), 即满足要求.\par
  由题目条件, \( \Lambda<1 \), 因此 \( \lim_{k\to\infty}k^n\Lambda^k = 0 \).  因此, 存在 \( m \) 使得 \( m^n\Lambda^m \le \varepsilon \). 从而对于任意的 \( v\in V \), 都有
  \[ \abs{T^mv}\le\varepsilon\abs{v}. \]
\end{proof}



\begin{problem}[6C-9]
设 $V$ 是有限维的. 设 $P \in \mathcal{L}(V)$ 使得 $P^2 = P$ 且 $\ker P$ 中任一向量都正交于 $\im  P$ 中任一向量. 证明: 存在 $V$ 的子空间 $U$ 使得 $P = P_U$.
\end{problem}

\begin{proof}
  记 \( U=\im P \). 则 \( U \) 是 \( V \) 的子空间. 下面来证明 \( P=P_U \).\par
  由于 \( \dim\ker P+\dim\im P=\dim V \), 且 \( \ker P \) 中任一向量都正交于 \( \im P \) 中任一向量, 因此 \( V=\ker P\oplus\im P \), 且 \( \ker P=U^\perp \). \par
  对于 \( u\in U=\im P \), 存在 \( v\in V \) 使得 \( u=Pv \). 由于 \( P^2=P \), 有 \( Pu=P^2v=Pv=u \). 对于 \( w\in U^\perp=\ker P \), 有 \( Pw=0 \). 因此对于任意的 \( v\in V \), 将 \( v \) 写成 \( v=u+w \) 的形式, 其中 \( u\in U \) 而 \( w\in U^\perp \), 则
  \[ Pv=Pu+Pw=u. \]\par
  而 \( u=P_U(v) \), 因此 \( P=P_U \).
\end{proof}



\begin{problem}[6C-10]
设 $V$ 是有限维的, $P \in \mathcal{L}(V)$ 使得 $P^2 = P$ 且
\[
\lVert Pv \rVert \le \lVert v \rVert
\]
对任一 $v \in V$ 成立. 证明存在 $V$ 的子空间 $U$ 使得 $P = P_U$.
\end{problem}

\begin{proof}
  先来证明: \( V=\ker P\oplus\im P \). 由于 \( \dim\ker P+\dim\im P=\dim V \), 只需要证明 \( \ker P\cap\im P=\{0\} \). \par
  取 \( u\in\ker P\cap\im P \). 由 \( u\in\im P \) 知存在 \( v\in V \) 使得 \( u=Pv \). 由 \( u\in\ker P \) 知 \( 0=Pu=P^2v=Pv=u \). 因此 \( u=0 \), 从而 \( \ker P\cap\im P=\{0\} \).\par
  记 \( U=\im P \). 与上一题类似地, 对于任意的 \( u\in U \), 均有 \( Pu=u \). 接下来证明: \( \ker P=U^\perp \). 由 \( \dim\ker P=\dim V-\dim\im P=\dim U^\perp \), 只需要证明 \( \ker P\subseteq U^\perp \). \par
  否则, 存在 \( v\in\ker P \) 且 \( v\not\in U^\perp \), 故 \( u:=P_U(v)\neq 0 \). 待定正实数 \( \lambda \), 并且考虑向量 \( v'=v-\lambda u \). 由题目条件, 有 
  \[ \abs{v'}\ge\abs{P(v')}=\abs{P(v)-\lambda P(u)}=\lambda\abs{u}. \]\par
  而
  \begin{align*}
    \abs{v'}^2
    &=\inprod{v-\lambda u}{v-\lambda u}
    =\abs{v}^2-2\lambda\Re\inprod{v}{u}+\lambda^2\abs{u}^2,
  \end{align*}
  因此, 对于任意的正实数 \( \lambda \), 都有
  \[ \abs{v}^2-2\lambda\Re \inprod{v}{u}+\lambda^2\abs{u}^2\ge \lambda^2\abs{u}^2\implies \abs{v}^2\ge 2\lambda\Re \inprod{v}{u}. \]
  而 \( \inprod{v}{u}=\inprod{u+(v-P_U(v))}{u}=\abs{u}^2>0 \), 因此上面的不等式不可能成立. 从而 \( \ker P\subseteq U^\perp \).\par
  因此 \( V=\ker P\oplus\im P \) 且 \( \ker P=U^\perp \). 因此 \( P=P_U \).
\end{proof}



\begin{problem}[6C-12]
设 $V$ 是有限维的, $T \in \mathcal{L}(V)$, 且 $U$ 是 $V$ 的子空间. 证明:
\[
U \text{ 和 } U^{\perp} \text{ 都在 } T \text{ 下不变 } \Longleftrightarrow P_U T = T P_U.
\]
\end{problem}

\begin{proof}
  ``\( \implies \)''. 对于任意的 \( v\in V \), 都可以做分解: \( v=P_U(v)+P_{U^\perp}(v) \). 则由条件, \( TP_U(v)\in U \), \( TP_{U^\perp}(v)\in U^\perp \). 因此
  \begin{align*}
    P_UT(v)
    &=P_UTP_U(v)+P_UTP_{U^\perp}(v)
    =TP_U(v).
  \end{align*}\par
  ``\( \impliedby \)''. 对于任意的 \( v\in V \), 有
  \[ P_UTP_{U^\perp}(v)=TP_UP_{U^\perp}(v)=T(0)=0,\quad
  P_{U^\perp}TP_U(v)=P_{U^\perp}P_UT(v)=0. \]
  因此, 对于 \( u\in U \), 有 \( P_{U^\perp}Tu=P_{U^\perp}TP_U(u)=0 \), 因此 \( Tu\in U \). 同理, 对于任意的 \( w\in U^\perp \), 有 \( Tw\in U^\perp \).\par
  因此 \( U \) 和 \( U^\perp \) 都在 \( T \) 下不变.
\end{proof}



\begin{problem}[6C-20]
设 $V$ 是有限维的, 且 $T \in \mathcal{L}(V, W)$. 证明:
\[
\ker T^{\dagger} = (\im  T)^{\perp}\quad \text{且}\quad
\im  T^{\dagger} = (\ker T)^{\perp}.
\]
\end{problem}

\begin{proof}
  由定义, 
  \[ T^\dagger=(T|_{(\ker T)^\perp})^{-1}\circ P_{\im T}. \]
  由于 \( (T|_{(\ker T)^\perp})^{-1} \) 是可逆的, 因此 \( \ker T^\dagger=\ker P_{\im T}=(\im T)^\perp \).\par
  另一方面, 
  \[ \im T^\dagger=\left( T|_{(\ker T)^\perp} \right)^{-1}P_{\im T}(W)
  =\left( T|_{(\ker T)^\perp} \right)^{-1}(\im T)= (\ker T)^\perp. \]
\end{proof}



\begin{problem}[6C-22]
设 $V$ 是有限维的, 且 $T \in \mathcal{L}(V, W)$. 证明:
\[
T T^{\dagger} T = T\quad \text{且}\quad T^{\dagger} T T^{\dagger} = T^{\dagger}.
\]
\end{problem}

\begin{proof}
  由定义, 对于任意的 \( v\in V \), 都有
  \begin{align*}
    TT^\dagger T(v)
    &=T\left( T|_{(\ker T)^\perp}\right)^{-1}P_{\im T} T(v)
    =T\left( T|_{(\ker T)^\perp}\right)^{-1}T(v)\\
    &=T(P_{(\ker T)^\perp}(v))=T(v)+T(P_{\ker T}(v))=T(v).
  \end{align*}\par
  对于任意的 \( w\in W \), 有
  \begin{align*}
    T^\dagger T T^\dagger(w)
    &=T^\dagger T\left( T|_{(\ker T)^\perp}\right)^{-1}P_{\im T}(w)
    =T^\dagger P_{\im T}(w)\\
    &=T^\dagger(w).
  \end{align*}\par
  因此, \( TT^\dagger T = T \) 且 \( T^\dagger T T^\dagger = T^\dagger \).
\end{proof}



\begin{problem}[6C-23]
设 $V$ 和 $W$ 是有限维的, 且 $T \in \mathcal{L}(V, W)$. 证明:
\[
(T^{\dagger})^{\dagger} = T.
\]
\end{problem}

\begin{proof}
  由前面的问题, \( \ker T^{\dagger} = (\im  T)^{\perp}\) 且 \( \im  T^{\dagger} = (\ker T)^{\perp} \). 因此 
  \[ T^\dagger|_{(\ker T^\dagger)^\perp}=T^\dagger|_{\im T}=\left( T|_{(\ker T)^\perp} \right)^{-1}
  \implies \left( T^\dagger|_{(\ker T)^{\perp}} \right)^{-1}=T|_{(\ker T)^\perp}. \]
  故
  \[ (T^\dagger)^\dagger = \left( T^\dagger|_{(\ker T)^{\perp}} \right)^{-1}P_{(\im T^\dagger)}
  =T|_{(\ker T)^\perp}P_{(\ker T)^\dagger}=T. \]
\end{proof}



\begin{problem}[7A-7]
证明: 如果 $T \in \mathcal{L}(V, W)$, 那么
\textbf{(a)} $\dim \ker T^* = \dim \ker T + \dim W - \dim V$;
\textbf{(b)} $\dim \im T^* = \dim \im T$.
\end{problem}

\begin{proof}
  (a)
  \[ \dim\ker T^*=\dim(\im T)^\perp = \dim W - \dim \im T= \dim W-\dim V +\dim \ker T. \]
  (b)
  \[ \dim\im T^*=\dim(\ker T)^\perp = \dim V - \dim \ker T = \dim \im T. \]
\end{proof}



\begin{problem}[7A-20]
设 $P \in \mathcal{L}(V)$ 使得 $P^2 = P$. 证明下列等价:
\textbf{(a)} $P$ 是自伴的.
\textbf{(b)} $P$ 是正规的.
\textbf{(c)} 存在 $V$ 的子空间 $U$ 使得 $P = P_U$.
\end{problem}

\begin{proof}
  (a) \( \implies \) (b) 是显然的. (c) \( \implies \) (a) 可以直接验证. 下面来证明 (b) \( \implies \) (c). \par
  设 \( U=\im P \). 由于 \( P \) 正规, 有 \( \ker P=\ker P^* \). 对任意 \( u\in\ker P \) 和 \( v\in U \), 存在 \( w\in V \) 使得 \( v=Pw \), 因而
  \[ \inprod{v}{u}=\inprod{Pw}{u}=\inprod{w}{P^*u}=0. \]
  故 \( U\perp\ker P \). 熟知 \( \dim U+\dim\ker P=\dim V \), 从而
  \[ V=U\oplus\ker P,\quad \ker P=U^\perp. \]\par
  对于 \( u\in U \), 有 \( Pu=P^2w=Pw=u \), 对于 \( w\in U^\perp=\ker P \), 有 \( Pw=0 \). 因此对任意 \( v=u+w \) (其中 \( u\in U \), \( w\in U^\perp \)),
  \[ Pv=Pu+Pw=u=P_U(v). \]
  从而 \( P=P_U \).
\end{proof}



\begin{problem}[7A-23]
设 $T$ 是 $V$ 上的正规算子. 设 $v, w \in V$ 满足下列各式:
\[\lVert v \rVert = \lVert w \rVert = 2, \quad Tv = 3v, \quad Tw = 4w.\]
证明: $\lVert T(v + w) \rVert = 10$.
\end{problem}

\begin{proof}
  \[ \inprod{T(v+w)}{T(v+w)}
    =\inprod{3v}{3v}+\inprod{4w}{4w}+2\Re\inprod{Tv}{Tw}
    =100+24\Re\inprod{v}{w}. \]\par
  由于正规算子的不同特征值对应的特征向量正交, 因此 \( \inprod{v}{w}=0 \). 从而 \( \lVert T(v+w) \rVert=10 \).
\end{proof}



\begin{problem}[7A-27]
设 $T \in \mathcal{L}(V)$ 是正规的. 证明:
\[\ker T^k = \ker T \quad \text{和} \quad \im T^k = \im T\]
对任一正整数 $k$ 都成立.
\end{problem}

\begin{proof}
  对 \( k \) 归纳. \( k=1 \) 时, 结论显然成立. 假设 \( k-1 \) 时结论成立. 下面来证明 \( k \) 时结论也成立.\par
  先来证明 \( \ker T^k=\ker T \). 由于 \( \ker T\subseteq\ker T^k \), 只需要证明 \( \ker T^k\subseteq\ker T \). 取 \( v\in\ker T^k \). 则 \( 0=T^k(v)=T(T^{k-1}(v)) \), 故 \( T^{k-1}(v)\in\ker T \). 显然 \( T^{k-1}(v)\in\im T \). 而由正规算子的性质, \( \ker T\cap\im T=\{0\} \), 因此 \( T^{k-1}(v)=0 \). 由归纳假设, \( v\in\ker T \).\par
  再来证明 \( \im T^k=\im T \). 注意到, 由于 \( T \) 正规, 故 \( T \) 与 \( T^* \) 可交换, 因此 \( T^k \) 的伴随为 \( (T^k)^*=(T^*)^k \). 因此 \( \im T^k=(\ker (T^*)^k)^\perp=(\ker T^*)^\perp=\im T \) (其中对 \( T^* \) 使用了上面的结论). 
\end{proof}



\begin{problem}[7A-29]
证明或给出一反例: 如果 $T \in \mathcal{L}(V)$ 且 $V$ 有一规范正交基 $e_1, \ldots, e_n$ 使得
\[\lVert Te_k \rVert = \lVert T^* e_k \rVert\]
对任一 $k = 1, \ldots, n$ 都成立, 那么 $T$ 是正规的.
\end{problem}

\begin{solution}
  反例: 设 \( V=\C^2 \). 取 \( T \) 的表示矩阵为
  \[ M=\begin{pmatrix}
    1 & 1+\i\\
    -1-\i & 2
  \end{pmatrix}. \]
  则 \( T^* \) 的表示矩阵 \( M^* \) 为 \( M \) 的共轭转置: 
  \[ M^*=\begin{pmatrix}
    1 & -1+\i\\
    1-\i & 2
  \end{pmatrix}. \]
  不难验证 \( \abs{Te_1}=\abs{T^*e_1}=\sqrt{3} \), \( \abs{Te_2}=\abs{T^*e_2}=\sqrt{6} \). 但 \( TT^*-T^*T \) 的表示矩阵为
  \[ MM^*-M^*M=\begin{pmatrix}
    0 & 2\\
    2 & 0
  \end{pmatrix}\neq 0. \]
\end{solution}



\begin{problem}[7A-32]
设 $T: V \to W$ 是一线性映射. 证明: 在 $V$ 和 $V'$ 的标准等同关系 (见 6.58) 以及相应的 $W$ 和 $W'$ 的等同关系下, 伴随映射 $T^*: W \to V$ 对应于对偶映射 $T': W' \to V'$. 更准确地说, 证明:
\[T'(\varphi_w) = \varphi_{T^* w}\]
对所有 $w \in W$ 成立, 其中 $\varphi_w$ 和 $\varphi_{T^* w}$ 如 6.58 所定义.
\end{problem}

\begin{proof}
  按照对偶映射的定义, 对于任意的 \( v\in V \), 
  \[ T'(\varphi_w)(v) = \varphi_w(Tv) 
    = \inprod{Tv}{w} 
    = \inprod{v}{T^*w} 
    = \varphi_{T^*w}(v). \]
\end{proof}



\begin{problem}[7B-6]
设 $V$ 是复内积空间, $T \in \mathcal{L}(V)$ 是使得 $T^9 = T^8$ 的正规算子. 证明: $T$ 是自伴的且 $T^2 = T$.
\end{problem}

\begin{proof}
  由复谱定理, \( T \) 关于某个单位正交基有对角矩阵. 而 \( T^9=T^8 \) 则说明 \( T \) 的每个特征值 \( \lambda \) 满足 \( \lambda^9=\lambda^8 \), 从而 \( \lambda=0 \) 或 \( 1 \). 因此 \( T \) 在这个单位正交基下的矩阵是一个以 \( 0 \) 和 \( 1 \) 为元素的对角矩阵, 从而 \( T \) 是自伴的且 \( T^2=T \).
\end{proof}



\begin{problem}[7B-10]
设 $V$ 是复内积空间. 证明: $V$ 上每个正规算子都有平方根.
\par
注: 算子 $S \in \mathcal{L}(V)$ 称为 $T \in \mathcal{L}(V)$ 的平方根 (square root), 如果 $S^2 = T$.
\end{problem}

\begin{proof}
  由复谱定理, \( T \) 关于某个单位正交基有对角矩阵. 设 \( \lambda_1, \ldots, \lambda_n \) 是 \( T \) 的特征值. 则对于每个 \( i \), 取 \( \mu_i \) 是 \( \lambda_i \) 的一个平方根. 则在这个单位正交基下, \( T \) 的矩阵是
  \[ M=\begin{pmatrix}
    \lambda_1 & &\\
    &\ddots&\\
    & &\lambda_n
  \end{pmatrix}, \]
  因此 \( S=\operatorname{diag}(\mu_1, \ldots, \mu_n) \) 满足 \( S^2=T \).
\end{proof}



\begin{problem}[7B-19]
设 $T \in \mathcal{L}(V)$ 是自伴的, $U$ 是 $V$ 的在 $T$ 下不变的子空间.
\textbf{(a)} 证明: $U^\perp$ 在 $T$ 下不变.
\textbf{(b)} 证明: $T|_U \in \mathcal{L}(U)$ 是自伴的.
\textbf{(c)} 证明: $T|_{U^\perp} \in \mathcal{L}(U^\perp)$ 是自伴的.
\end{problem}

\begin{proof}
  (a) 对于任意的 \( w\in U^\perp \), 我们要证明 \( Tw\in U^\perp \). 而
  \[ Tw\in U^\perp\Leftrightarrow \inprod{Tw}{u}=0,\forall u\in U\Leftrightarrow \inprod{w}{Tu}=0,\forall u\in U. \]
  由 \( w\in U^\perp \) 以及 \( Tu\in U \) 知 \( \inprod{w}{Tu}=0 \), 因此 \( Tw\in U^\perp \).\par
  (b) 对于任意的 \( u_1, u_2\in U \), 有
  \[ \inprod{T|_U(u_1)}{u_2}=\inprod{T(u_1)}{u_2}=\inprod{u_1}{T(u_2)}=\inprod{u_1}{T|_U(u_2)}. \]
  因此 \( T|_U \) 是自伴的.\par
  (c) 结合 (a) 与 (b), \( T|_{U^\perp} \) 是 \( U^\perp \) 上的自伴算子.
\end{proof}



\end{document}